\documentclass[../main.tex]{subfiles}
\begin{document}
%==========================================TIÊU ĐỀ TẬP========================================
\begin{table}[h]
    \begin{adjustwidth}{-1cm}{}
        \begin{tabular}{>{\centering\arraybackslash}p{6cm}|>{\raggedright\arraybackslash}p{12.5cm}}
            \multirow{5}{6cm}
            % Cột bên trái: ÉPISODE và số tập
            {\\[-40pt]
                \flaregothic\fontsize{20pt}{20pt}\selectfont \centering
                \color{eptype}HISTOIRE\color{black}\\
                \freshbot\fontsize{72pt}{72pt}\selectfont
                \color{epnum}4
                \\[18pt]
                \flaregothic\fontsize{14pt}{14pt}\selectfont
                \color{epattr}BIENVENUE À NOUVEAU, \uline{\color{Mtfk} Coco} !
                % \flaregothic\fontsize{18pt}{18pt}\selectfont
                % \color{epattr}ÉPISODE PERDU
                \color{black}
            }
             &
            % Dòng thứ nhất: tên tập tiếng Nhật
            {\fotskip\fontsize{18pt}{18pt}\selectfont おはようございます！}
            \\[6pt]
            % Dòng thứ hai: tên tập tiếng Anh
             & {\cafeteria\fontsize{18pt}{18pt}\selectfont Good Afternoon, Motherf-kers!}            \\[4pt]
            % Dòng thứ ba: tên tập tiếng Việt
             & {\vnmsans\fontsize{18pt}{18pt}\selectfont Chào buổi chiều, những kẻ ở xứ Kawachi!}    \\[8pt]
            % Dòng thứ tư: Nhân vật xuất hiện
             & {\brian{\fontsize{14pt}{14pt}\selectfont Track used: Jurassic Marsh - Ultimate Battle}
               }
            \\[8pt]
            % Dòng cuối cùng: Thời gian diễn ra của tập (thường sẽ là ngày phát hành)
             & {\continuum\fontsize{16pt}{16pt}\selectfont Ngày phát hành: 27/06/2021}               \\[18pt]
        \end{tabular}
    \end{adjustwidth}
\end{table}
%=========================================TRUYỆN CHÍNH========================================
\color{black}

\begin{center}
    (Tiếp tục từ phần {\abv Truyện phụ} của {\flaregothic ÉPISODE 112}, quyển số Chín.)
\end{center}

\pov{Hikawa Kuon}

Tôi có nói rằng tôi sẽ không bao giờ gặp lại Coco-chan nữa, đúng không nhỉ? Nhất là sau cái kiểu chào nhà binh đó - khi em ấy hóa thành một con rồng khổng lồ và bay đi như một cảnh trong phim hành động Hollywood chứ?

\dots Đúng là tôi đã nghĩ vậy.

Nhưng mà, tôi đã nhầm.

Có lẽ do tôi quá mải làm việc, đến mức quên mất một điều quan trọng: căn nhà của chúng tôi bây giờ đang là ``ký túc xá dành cho những thành viên đã tốt nghiệp'' theo yêu cầu và mong muốn của YAGOO.

Và người đầu tiên chuyển vào chính là Mano Aloe-chan hồi Tết Nguyên Đán vừa rồi.

Mấy hôm nay tôi không gặp Aloe-chan nhiều, vì em ấy bận bịu đi hát với đi làm. Nhưng thi thoảng, cô nàng Succubus ấy vẫn sẽ rủ tôi ra ngoài, hoặc nếu không thì lại kéo tôi xuống phòng em ấy để chơi điện tử, hát hò. Về cơ bản, vẫn là một con nhóc kusogaki hỗn hào như thường ngày thôi, không có gì khác lạ cả.

À, mà nhân tiện nói luôn, Tanik- à nhầm, Tanigo-san - ông sếp của tôi ấy, đã có ý định xây thêm vài ba tầng nữa trên nóc nhà.

``Đề phòng lại có thêm vài thành viên `không may' phải về hưu nữa, cậu cứ xây cho tôi thêm vài tầng nữa đi.''

Đó là một khả năng chẳng ai muốn nghĩ tới, nhưng\dots\ dù nghe có hơi điên rồ, tôi phải công nhận rằng ông ấy luôn chuẩn bị kỹ lưỡng.

Mặc dù vậy, tôi thực sự hy vọng rằng nhà chúng tôi sẽ không phải sử dụng mấy cái tầng đó\dots

\dots Nhưng không.

Có vẻ như Aloe-chan sẽ sớm có bạn cùng ký túc xá đây.

\il

Buổi chiều ngày hôm đó.

Sau cái chào nhà binh ``lịch sử'' ấy, cả Towa-chan, Kanata-chan, Luna-chan và Watame-chan được công ty cho nghỉ sớm, để chuẩn bị cho buổi tốt nghiệp chính thức cho Coco-chan trong những ngày sắp tới.

Còn tôi? Tôi vốn dĩ định làm việc đến nát đít như mọi ngày.

Nhưng rồi, ông sếp của tôi - Tanigo-san - gọi tôi lên văn phòng và nói một câu ngắn gọn: ``Hôm nay về nhà sớm đi. Các cậu sẽ bàn giao người mới cho ký túc xá.''

Người mới?

Các cậu?

Ông ấy đang chỉ đến gia đình chúng tôi đúng không?

Đúng là như vậy. Vì khi tôi và bố tôi không có nhà, thì bác Hijaki-san sẽ là người giám sát.

Mặc dù bác ấy cũng không quản lý quá chặt, vì dù sao cũng đã hơn sáu mươi tuổi rồi, nhưng cuối tuần nào bác cũng báo cáo tình hình lại cho chúng tôi. Và hôm nay chính là cuối tuần.

Nói đến đây, có lẽ các bạn cũng muốn biết dạo này Aloe-chan thế nào rồi nhỉ?

Theo như lời bác Hijaki-san kể lại (và cũng từ mấy lần hiếm hoi tôi ghé thăm phòng con bé), Aloe dạo này đã có một công việc ổn định hơn.

Ngoài việc đi làm thêm ở hội chợ với mẹ tôi để lấy kinh nghiệm, em ấy vẫn tiếp tục phụ giúp ở quán café gần nhà, đồng thời chăm sóc cho bà nội tôi.

Tóm lại, mọi chuyện vẫn diễn ra bình thường.

Chưa có chuyện gì quá đặc biệt xảy ra cả.

\dots Cho đến giờ phút này.

Bởi vì nhà Hikawa chúng tôi sắp có một thành viên đã tốt nghiệp nữa - người đã tự nguyện chọn nơi đây làm chỗ ở tiếp theo.

Và nếu các bạn đã theo dõi đến tận đây, thì chắc hẳn cũng đã biết người đó là ai rồi.

Chỉ nghĩ đến thôi, tôi đã thở dài, cảm nhận rõ ràng thứ mà mình sắp phải chạm trán tới\dots

\begin{center}
    \uline{Nhà Hikawa.}
\end{center}

Và, đúng như tôi dự đoán. Một mớ bòng bong đã chào đón tôi ngay khi tôi về đến nhà.

``\vie{\dots Nghiêm túc đấy à?}''

Ngay trước cửa nhà tôi là một đám đông tụ tập, có người giơ điện thoại lên quay, có người thì xì xào bàn tán, có người lại vừa cười vừa trầm trồ.

Bình thường thì tôi sẽ chẳng thèm quan tâm, nhưng linh cảm của tôi mách bảo rằng chuyện này có liên quan đến \hll.

Người Kawachi nói riêng và người Etsunan nói chung vốn hiếu kỳ từ trong máu.

Chỉ cần có một người cưa bom là cả hội sẽ kéo đến xem, không ai quan tâm hậu quả ra sao, chỉ cần vui là được. Vậy nên, việc họ đứng đông nghẹt trước cửa nhà tôi cũng chẳng có gì quá lạ lẫm.

Thậm chí, người dân ở đây còn nhiệt tình tụ tập đánh hội đồng đến nhừ tử một đám trộm chó có vũ trang nếu chẳng may chúng sa vào lưới của họ. Những chuyện như thế, tôi đã từng nghe qua.

Nhưng lần này thì khác.

Tôi cố chen lấn qua đám đông, len lỏi vào nhà, và ngay khi vừa bước vào tầng một, thứ chào đón tôi là\dots

Một bãi chiến trường.

Căn tầng một - vốn là nơi gia đình tôi cho thuê làm cửa hàng viễn thông - giờ tan hoang như vừa bị đánh bom.

Cửa kính vỡ tan tành.

Kệ tủ với bàn làm việc chi chít vết đạn.

Giấy tờ bay tán loạn.

Thậm chí một mảng tường còn bị thủng một lỗ to tổ bố, mà tôi cá là được bắn ra từ một khẩu bazooka hoặc một khẩu phóng lựu.

\dots Rất may là không có ai chết.

Pekora-chan hôm nay đến nhà tôi chơi à?

Mà khoan, cho dù có là nàng Thỏ chiến tranh, tôi cũng không nghĩ em ấy có thể quậy banh nóc đến mức này. Với lại, hôm nay em ấy có buổi phát sóng trực tiếp, thì càng không thể có mặt ở đây được.

Khi tôi còn đang cố sắp xếp lại suy nghĩ, chợt nghe thấy giọng của một nhân viên trong cửa hàng, đang bàn tán với người dân xung quanh.

``\vie{Có một cô gái điên rồ nã đạn vào cửa hàng chúng tôi!}'' nữ nhân viên ấy nói với mọi người khi một trong số những người trong đám đông tra hỏi.

\dots Ờm, chuyện này bắt đầu nghe có vẻ quen quen rồi đấy.

Khu tôi ở hiếm khi xảy ra những vụ cướp, mà lại còn do một cô gái thực hiện?

Lại còn thẳng tay tấn công cả một cửa hàng viễn thông giữa kinh thiên bạch nhật nữa?

Cảm thấy có gì đó bất ổn, tôi tiến tới hỏi người nhân viên đó.

``\vie{Mọi chuyện xảy ra thế nào ạ?}'' tôi hỏi nữ nhân viên đang hoảng loạn kia.

Cô ấy đáp, khuôn mặt vẫn còn xanh lét: ``\vie{Đây nhá, đang yên đang lành thì bỗng dưng có một cô gái cao từng này\dots}'' Người đó giơ cánhtay lên cao hơn tôi, ám chừng khoảng mét tám mươi.

``\vie{Cô ta cao đến thế ạ!?}'' tôi thoáng một chút ngạc nhiên, rồi tiếp tục tra hỏi: ``\vie{Rồi, gì nữa ạ?}''

Nữ nhân viên lắp bắp trả lời, cố nhớ lại sự việc kinh hoàng trên: ``\vie{Cô ta hét lớn:} `\eng{\textbf{Good afternoon, motherf*ker!!}}' \vie{rồi đột nhiên xả hàng loạt ổ đạn vào chúng tôi! Thậm chí còn nã súng phóng lựu vào tường nữa!}''

\dots

Sao tôi thấy câu chào này\dots

\dots có vẻ quen thuộc nhỉ?

Cảm thấy vẫn chưa hoàn toàn tin tưởng, tôi tiếp tục hỏi: ``\vie{Cô ta trông như thế nào ạ?}''

Nữ nhân viên ấy kể tường tận cả ngoại hình mà nữ sát thủ kia - tôi cho là vậy - diện: ``\vie{Tóc màu cam dài đến lưng, trên đầu có hai cái sừng, mặc bộ đồ giống nữ sinh trung học: áo khoác đen, đầm đỏ\dots}''

``\vie{Vâng\dots}''

``\vie{Ngực rất to, trên cổ đeo một cái dây chuyền có ngọc tím đỏ hình lửa, lại còn có một cái đuôi rồng màu tím thù lù ở phía sau nữa!}''

\dots Không cần nghe thêm nữa.

Tôi đã biết đó là ai rồi.

Dám chắc đây là người mà Tanigo-san bàn giao cho tôi đây mà.

\dots Tự nhiên lại nhớ đến hồi năm ngoái, khi em ấy lần đầu lên văn phòng và đòi chơi ``đột kích'' với chúng tôi.

Cũng y như vậy. Nhưng lần này em ấy chơi tới bến luôn, chứ không làng nhàng như hồi trước nữa.

Thở dài một hơi, tôi hỏi tiếp: ``\vie{Vậy cô ấy đang ở đâu ạ?}''

Người nhân viên vẫn còn run rẩy, chỉ tay hướng lên tầng năm, chỗ căn phòng của hai bố con tôi.

Tôi không chần chừ thêm giây nào nữa.

``\vie{N-Này cậu bé!?}'' / ``\vie{Cậu chạy đi đâu vậy!?}'' / ``\vie{Quay lại đây, đồ ngốc!}''

Ngay lập tức, tôi chạy thẳng vào trong nhà, mặc kệ những lời ngăn cản của người dân xung quanh.

\ct

Tôi lao nhanh lên cầu thang, lòng nóng như lửa đốt.

Lần lượt đi qua các phòng trên lầu, tôi vẫn không thấy dấu hiệu gì bất thường - ít nhất là cho đến lúc đó.

Phòng của Aloe? Cửa đóng then cài. Chắc em ấy đang đi làm cùng với mẹ tôi.

Phòng của Ryuuhou và mẹ tôi? Không có ai ở đó, vì em tôi đang sang chơi bên nhà bà ngoại ở ngoại thành, còn mẹ tôi thì vẫn đang đi làm.

Phòng của bà nội? Bà tôi đang ngủ say, và bác Hijaki-san cũng vậy.

\dots Đợi đã.

Tiếng súng nổ long trời lở đất như vậy\dots

Mà bác Hijaki-san vẫn ngủ ngon lành!?

Tôi liếc vào trong phòng, thấy bác ấy đang kê tay sau đầu, ngủ ngon như chưa từng có vụ nổ nào vừa xảy ra.

\dots Tôi không biết nên cảm thấy khâm phục hay lo ngại nữa.

Khi tôi lên đến phòng của bố con tôi, cánh cửa trượt đã mở toang.

Mi, con chó cảnh lông trắng của tôi, sủa ầm ĩ, mõm hướng thẳng về phía phòng tắm.

Có ai đó ở trong đó.

Tôi bước tới, tay đặt lên cánh cửa\dots

Nhưng bỗng nhiên\dots

*ĐOÀNG!*

Cánh cửa bật tung ra, gãy luôn cả bản lề, suýt nữa va thẳng vào mặt tôi.

Rất may là tôi có phản xạ nhanh, giật người sang một bên tránh theo bản năng.

Và trước mặt tôi, đúng như miêu tả của người nhân viên ban nãy:

Tóc cam dài đến lưng.

Đuôi rồng màu tím phe phẩy sau lưng.

Mặc đồng phục nữ sinh trung học: áo khoác đen, váy đỏ.

Trên hông giắt một khẩu bazooka và một khẩu súng trường.

Ngoài ra còn có đeo mặt nạ Jason từ ``\textit{Thứ Sáu ngày 13}''.

Trên tay còn đang cầm\dots\ một cái cưa máy.

Tôi cạn lời.

Trước khi tôi kịp mở miệng, người đó đã cất giọng trước. ``Ô, Kuon-paisen!''

Chậc. Một giọng nói trẻ con quen thuộc.

``Anh ở đây làm gì vậy?''

Tôi chớp mắt, day trán khẽ lắc đầu, rồi ngán ngẩm đáp lại: ``\dots Sao với trăng gì em. Đây là nhà riêng của anh mà.''

Không ai khác, đây chính là Kiryu Coco - người mà tôi đã chào binh ở cửa sổ vào sáng cùng ngày.

Nhìn thấy bóng dáng của cô nàng Rồng quậy phá quen thuộc, tôi chỉ có thể thở dài, rồi xoa trán: ``Em làm cái trò gì đấy, Coco-chan? Cái mặt nạ Jason với cái cưa máy là thế nào đây?''

``Em đang chơi trò `đột kích' đấy!'' em ấy tháo mặt nạ ra, hớn hở đáp.

\dots Chết tiệt, đoán không sai mà.

``Chả trách có cả súng ống giắt trên người kia kìa. Mà hình như em đang làm thật đến mức quá đáng rồi đấy.''

Coco chớp mắt, tỏ vẻ khó hiểu.

Tôi chỉ lắc đầu ngao ngán, chỉ tay về phía cửa sổ.

Coco tò mò nghiêng đầu, rồi làm theo lời tôi\dots

Và dưới nhà, cảnh tượng đã khác hẳn.

Đám đông trước cửa nhà giờ còn đông hơn. Hàng chục ống kính từ điện thoại giơ lên, hướng về phòng của bố con tôi.

Thậm chí có cả xe cảnh sát, bao vây kín khu vực. Lực lượng an ninh tới đây còn can thiệp để vừa giải tán đám đông, vừa có ý định xông vào điều tra đây.

Họ mà biết được có một cô Rồng làm loạn thế này thì có mà hỏng việc!

Tôi khoanh tay, nhìn những món hàng nóng mà em ấy mang theo: ``Anh thấy em đang giả lập làm tội phạm thì đúng hơn đấy. Mà em làm mấy thứ ba lăng nhăng này để làm cái quẹc gì?''

``Ể? Muốn có một màn chào hàng mừng em trở thành một thành viên của ký túc xá này mà cũng không được sao?'' Coco-chan nhíu mày, giở giọng hơi thất vọng.

Tới lúc này, tôi chẳng biết em ấy đang nghĩ gì nữa. Biết là em ấy trẻ con, nhưng đến mức trẻ trâu tới mức này thì tôi cũng muốn vái cả nón.

Tôi hít một hơi, rồi giơ ngón trỏ lên.

``Thứ nhất, em vẫn còn một buổi tốt nghiệp hẳn hoi nữa trong vài ngày tới, thế nên em chỉ đến đây để bàn giao bước đầu thôi, chứ chưa phải là thành viên chính thức đâu.''

Tôi giơ tiếp ngón giữa, tiến gần một bước nữa.

``Thứ hai, nếu mà chào hàng kiểu này thì thôi, thà em thuê chỗ nào khác mà ở chứ làm ơn đừng quấy rối ở đây. Kawachi chứ không phải như thành phố Điện tử đâu.''

Tôi giơ tiếp ngón áp út, dí gần hơn nữa.

``Và thứ ba, nghe không liên quan, nhưng mà nhiều vũ khí trên người thế không sợ nặng à?''

Mỗi điều tôi nói, tôi tiến dần về phía Coco-chan, làm cho em ấy phải lùi dần về phía sau.

Trong khi đó, Mi vẫn tiếp tục sủa oăng oẳng như đồng tình với ý kiến tôi.

Tôi quắc mắt ``Cả mày nữa, ngậm cái mồm vào!''

Con chó giật mình, rồi vẫn tiếp tục sủa, như thể nó coi đây là nghĩa vụ thiêng liêng.

Còn Coco-chan?

Em ấy chỉ cười khẽ, rồi lé lưỡi ra, tỏ vẻ ``ngây thơ vô số tội''.

Chúng tôi liếc xuống dưới. Trước cửa nhà giờ đây còn đông đúc như một đám quân Nguyên, với hàng loạt máy quay, thiết bị bay, rồi còn cả các thiết bị chuyên dụng để tác nghiệp nữa!

``Bây giờ làm thế nào thì làm, giải tán cái đám người ở dưới nhà hộ anh cái,'' tôi gằn giọng, vừa thở dài vừa xoa thái dương.

Coco-chan không nói không rằng, tiến ngay đến cửa sổ, móc trong túi ra mấy viên gì đó màu trắng rồi ném thẳng xuống đám đông bên dưới. Trong chớp mắt, một làn khói trắng đục bốc lên mù mịt, che khuất cả tầm nhìn.

Tôi giật mình, nhướng mày hỏi: ``Em vừa làm cái quái gì thế?''

Không thèm trả lời, em ấy tiếp tục ném thêm vài viên nữa, khiến cho lớp khói càng dày đặc hơn. Giờ thì quang cảnh trông chẳng khác gì buổi sáng sớm ban mai vào mùa đông với một bầu trời đầy sương mù cả!

Nhưng, điều mà khi khói bắt đầu tan dần làm cho tôi sốc đến nhức óc.

Tôi có thể thấy rõ phản ứng của đám đông: người thì đứng đờ ra, gãi đầu gãi tai như thể quên mất mình đang làm gì, người thì cười ngơ ngẩn như thể vừa nghe chuyện cười hay nhất thế giới, người thì rú lên rồi chạy tán loạn không lý do.

Chứng kiến cảnh tượng ấy, tôi đổ mồ hôi hột: ``Chuyện quái gì đang xảy ra vậy?''

Coco nở nụ cười đầy tự hào: ``Nhờ AsaCoco của em đấy!''

Tôi nheo mắt, nghiêng đầu đầy khó hiểu: ``AsaCoco? Cái đó là cái gì?''

Nàng Rồng hào hứng đáp lại, kể ra một tràng mà tôi nghe còn không hiểu: ``Một loại hỗn hợp đặc chế từ tay em! Vừa có tác dụng kích thích, vừa có tác dụng giảm đau, đảm bảo khiến người dùng rơi vào trạng thái ảo giác siêu thực!''

Chất quái gì lạ vậy? Tôi chưa bao giờ nghe đến loại hóa chất nào như thế.

Tôi cau mày, nheo mắt hỏi cô nàng: ``Nói trắng ra thì\dots\ nó là thuốc gây ảo giác?''

Thậm chí nếu tệ hơn, thì đó là một dạng như thuốc phiện ma túy. Một thứ chất cấm mà nếu cảnh sát biết được thì chỉ có đường chết.

``Có thể hiểu như thế ha?'' Coco-chan gật đầu, vẻ đắc ý. ``Em còn làm cả một series quảng cáo về nó nữa cơ!''

Tôi nhìn em ấy với ánh mắt không thể tin nổi. Quảng cáo cho cái chất bột trắng trông như chất cấm này á!?

``Anh rất ngạc nhiên khi em nói mấy thứ này một cách thản nhiên đấy\dots\ Với cả, ai lại quảng cáo cái thứ-'' Tôi bỗng im bặt khi nhìn ra ngoài cửa sổ lần nữa.

Cái gì đây\dots?

Trước cửa nhà tôi, nơi từng có cả đống người vây quanh chỉ vài phút trước, giờ đã hoàn toàn vắng tanh. Không một bóng người, không một chiếc xe, thậm chí ngay cả đường phố cũng trông có vẻ hoang vắng hơn bình thường. Một cảm giác rờn rợn chạy dọc sống lưng tôi.

``T-Thật hả giời\dots'' tôi lẩm bẩm nói.

Coco vỗ vai tôi một cách đầy tự hào: ``Thấy chưa? Em có thể xử lý mọi thứ như một chủ bang yakuza chính hiệu luôn đó nha!''

Tôi lườm em ấy: ``Rồi rồi, anh ghi nhận. Nhưng mà là, ơn bớt cái giọng khoe khoang lại đi, Kiryu Yakuza ạ. Chưa bị tống vào nhà đá bóc lịch là may lắm rồi.''

Cô nàng cười khoái chí, vẫy chiếc đuôi to thù lù ở đằng sau: ``Chuyện nhỏ! Cái này em làm quen rồi!''

Tôi đổ mồ hôi hột: ``Anh không có khen em đâu mà cười tươi rói thế.''

Chợt nhớ đến cái cửa phòng tắm vừa bị đá bay ban nãy, tôi thở dài, chỉ tay về phía nó: ``Và bây giờ thì, phiền em đền cho nhà anh cái cửa nhà tắm này nha, bọn anh cần gấp đấy.''

Coco-chan cười tự tin: ``Cái này thì dễ!''

``Và cả nguyên cái cửa hàng viễn thông ở tầng một nữa,'' Tôi tiếp lời, mặt không cảm xúc. ``Phá được thì sửa được, đúng chứ?''

Nàng Rồng nghe đến đây liền cứng họng trong vài giây, rồi chỉ biết gãi đầu cười gượng.

Tôi khoanh tay nhìn Coco với ánh mắt đầy nghi ngờ: ``Giờ thì em tính sao đây, Bang chủ?''

Nàng Rồng chống nạnh, mặt hớn hở như thể vừa lập được một kỳ tích vĩ đại: ``Anh cứ yên tâm! AsaCoco có thể giải quyết mọi vấn đề!''

``Tuyệt. Cụ thể là như thế nào?''

Coco-chan vỗ ngực tự tin, rồi không biết lôi từ đâu ra một cái lọ nhỏ chứa bột trắng. ``Chỉ cần rắc một ít AsaCoco phiên bản `Sửa Chữa Toàn Diện' lên chỗ hỏng, đảm bảo nó sẽ tự lành lại ngay lập tức!''

``Nghe khả nghi đấy\dots'' tôi lùi lại vài bước.

Cô nàng dứt khoát đổ thẳng một lượng lớn bột trắng lên khung cửa đã gãy. Trong tích tắc, có một tiếng *PÓP* vang lên, và\dots

Bằng một cách nào đó, cửa phòng tắm của tôi đã biến mất hoàn toàn.

Không phải là sửa chữa hay lành lặn gì cả. Mà là không còn cái cửa nào luôn.

Tôi ôm trán, còn Coco thì đứng trơ ra một lúc rồi gãi đầu: ``Ờm\dots\ chắc là em hơi lạm dụng liều lượng.''

Tôi cười khẩy: ``Vậy chắc cũng có AsaCoco phiên bản `Triệu Hồi Cửa Mới' chứ?''

``À\dots\ cái đó thì\dots'' Coco đảo mắt, nhìn quanh như đang tìm đường chuồn.

``Không có chứ gì?'' Tôi nhướng mày.

``\dots Có, nhưng em có thể-''

``Có thể gì? Để anh đoán nhé, em sẽ gợi ý anh uống AsaCoco rồi tự tưởng tượng là có cửa mới đúng không?''

Coco bật ngón tay cái, mặt đầy tự hào: ``Anh hiểu nhanh đấy, Kuon-paisen!''

Tôi thở dài, lắc đầu ngao ngán: ``Rồi, bỏ đi. Giờ thì phiền em đi mua cái cửa mới hộ anh.''

Nàng Rồng bĩu môi ngạc nhiên: ``Bắt em vác cái cửa mới về á? Nặng lắm đấy nhé!''

``Ờ, thì cũng như cái đống súng ống em vác trên người thôi. Mà nếu không thích thì đơn giản lắm\dots'' Tôi nở một nụ cười đầy ẩn ý, giơ ngón trỏ lên. ``Thứ nhất, nếu em không chịu đền, thì anh sẽ nói chuyện với Hijaki-san xem `người mới' ở ký túc xá định phá banh nhà thế nào.''

Coco-chan tái mặt.

Tôi giơ ngón giữa. ``Thứ hai, bố anh mà biết thì chắc chắn em sẽ phải nghe nguyên bài thuyết giảng dài bốn tiếng đồng hồ.''

Coco nuốt nước bọt. Tôi dọa thế mà em ấy cũng tin sao?

\dots Mà, có lẽ ông ấy sẽ mắng xa xả chúng tôi thật.

Rồi, tôi giơ ngón áp út. ``Thứ ba, nếu em không thích tự vác cửa về thì cứ thử dùng AsaCoco đi. Biết đâu em có thể khiến cửa tự mọc lại thì sao?''

Tới đây, cô nàng rùng mình: ``Ôi trời, em thà đi vác cửa còn hơn nghe bác trai giảng đạo!''

``Anh cũng nghĩ thế,'' tôi khoanh tay, cười một cáchthỏa mãn.

Rốt cuộc, cả tôi và em ấy phải trích tiền lương để đền bù nguyên một cánh cửa nhà tắm, một cánh cửa hàng viễn thông và\dots\ có lẽ là cả một khu phố vừa tạm thời trở thành ``vùng đất ma'' vì AsaCoco - mà ơn giời tôi không phải trả, vì thứ chất bột này sẽ hết tác dụng sau một khoảng thời gian, nhưng những gì xảy ra hôm nay là mọi người vẫn sẽ quên hết.

Màn chào hàng của em ấy kết thúc bằng một hóa đơn thanh toán dài ngoằng, và một bài học quan trọng: đừng bao giờ để một cô Rồng giải quyết vấn đề bằng thứ gọi AsaCoco. Chào hàng kiểu này thì có mà đuổi người ta chạy mất dép chứ chẳng đùa.

\ct

Sau khi mọi chuyện ổn định, tôi và Coco phải mất cả buổi để xử lý hậu quả do AsaCoco gây ra. Ngoài việc thay cửa nhà tắm và đền bù tổn thất cho cửa hàng viễn thông ở tầng một, chúng tôi còn phải đảm bảo rằng không ai trong khu phố kiện tụng gì về ``làn khói trắng bí ẩn'' đã khiến họ hành xử như người mất trí.

Khi mọi thứ đã đâu vào đấy, tôi mới thực sự có thời gian suy nghĩ về giải pháp dài hạn. Nếu cứ tiếp tục sống trong một ngôi nhà có tỷ lệ bị phá hủy cao như thế này, sớm muộn gì tôi cũng phải tìm cách sửa chữa nhanh hơn là gọi thợ sửa mỗi lần. Một cái máy tái tạo kiến trúc (regenerator) có lẽ sẽ là một giải pháp không tồi - một thiết bị có thể khôi phục lại trạng thái nguyên vẹn của ngôi nhà chỉ trong vài phút, thay vì tốn hàng tuần để sửa chữa từng chút một. Có thể tôi cũng sẽ áp dụng nó cho cả công ty nữa.

Nhưng đó là chuyện của một ngày khác. Vì quan trọng nhất bây giờ là phải giải quyết vấn đề chỗ ở của Coco. Cụ thể hơn, tôi cần đưa em ấy đi gặp bố tôi và bác tôi để thương lượng về việc em ấy chính thức chuyển vào ký túc xá nhà Hikawa trong những ngày sắp tới.

\il

Bố tôi khoanh tay, nhìn tôi chằm chằm vị khách siêu quậy với vẻ mặt khó tin sau khi chúng tôi kể lại hành trình Coco-chan tới ký túc xá của chúng tôi.

``Vậy\dots\ để bố tóm tắt lại. Có một cô Rồng chạy đến nhà mình làm loạn dưới tầng một vì nghĩ rằng mình đang đóng giả 'làm đột kích', sau đó phá tung cửa nhà tắm, rồi ném thứ bột trắng xuống đường khiến cả phố bị rối loạn. Đúng không?''

Tôi thở dài: ``Chính xác là như vậy đấy, bố ạ.''

Bác Hijaki-san lúc này mới ngồi thẳng dậy, nheo mắt nhìn tôi như thể tôi vừa kể một câu chuyện hoang đường: ``Và người làm những việc đó trong khi tao đang ngủ\dots\ lại chính là vị khách thuê phòng tiếp theo được sếp mày gửi đến đúng chứ?''

Tôi liếc sang cô nàng Rồng, em ấy chỉ cười ngượng ngùng. Tôi quay lại nhìn bác tôi, gật đầu một cách đầy miễn cưỡng. ``\dots Đúng nốt ạ.''

Không khí trong phòng rơi vào khoảng lặng khó xử.

Tôi, bố tôi và bác tôi đang ngồi trong phòng của bác Hijaki-san để bàn bạc về việc đón tiếp thành viên mới của ký túc xá nhà Hikawa - Kiryu Coco. Sau khi nghe tôi và em ấy kể lại toàn bộ sự việc, cả hai người lớn chỉ biết thở dài, lắc đầu ngao ngán.

``Dù bố biết công ty nơi con làm có rất nhiều chuyện quái gở,'' bố chẹp miệng, ``nhưng đến mức làm náo loạn cả một khu phố thì bố chưa bao giờ thấy ai làm như thế đấy.''

Tôi cũng không thể phản bác, chỉ đáp lại một cách mệt mỏi: ``Nói thế này không tin chứ\dots\ lúc mà Coco-chan lần đầu đến công ty cũng bày trò đột kích kiểu này đấy.''

Bố tôi nhướng mày ngạc nhiên: ``Thế à?''

``Vâng. Lúc thì làm người đưa thư, lúc thì cầm cưa máy, rồi cả cosplay khủng long, thậm chí còn phá tung cửa sổ đến mức A-san phải mắng cho một trận. Nhưng đến độ quá quắt thế này thì đây là lần đầu ạ.''

Bác Hijaki-san thở dài lần nữa, nhưng lần này có phần tò mò hơn: ``Thế ngoài mấy trò điên rồ đó, nó còn làm gì khác nữa không?''

Tôi ngẫm nghĩ trong vài giây, rồi nhún vai: ``Nói thật với bác chứ\dots\ em ấy còn làm nhiều chuyện giời ơi đất hỡi nữa, nhưng mà mọi chuyện cũng giải quyết xong hết rồi, đúng chứ, Coco-chan?''

Coco nở nụ cười rạng rỡ, ngẩng đầu đầy tự hào: ``Như anh ấy nói đấy! Kuon-paisen chịu bỏ tiền túi ra để đền bù một phần thiệt hại sau khi làm bom nổ này, rồi-''

``Em không cần phải nhắc lại chuyện đó đâu,'' tôi giơ tay lên ra hiệu dừng, mặt nhăn lại, méo xệch. ``Đau ví lắm\dots''

Bác trưởng nhà hắng giọng một cái, khiến cả ba chúng tôi im lặng ngay lập tức. Bác đưa mắt nhìn Coco-chan, rồi hỏi bằng một giọng trầm tĩnh nhưng đầy uy lực: ``Thôi, mấy cái chuyện linh tinh này tí nói sau. Bây giờ, hẳn là mày đã gây ra tội trạng gì nên mới bị đày đến đây hả?''

Nàng Rồng hơi giật mình trước câu hỏi thẳng thắn ấy, nhưng nhanh chóng lấy lại bình tĩnh. Cô nàng đáp gọn với sự tự tin: ``Dạ không đâu bác ạ, cháu tự nguyện xin tốt nghiệp mà bác.''

Bố tôi chống cằm, ánh mắt thoáng vẻ ngờ vực: ``Cái này thì hơi lạ nha. Cháu đang có công việc ổn định như thế mà lại rút sớm thế? Chú cũng thấy ngạc nhiên đấy.''

Coco-chan cúi đầu một chút, có vẻ trầm tư, nhưng ánh mắt vẫn giữ được nét kiên định: ``Thật ra\dots\ cháu đã suy nghĩ về chuyện tốt nghiệp từ lâu rồi ạ. Cháu cảm thấy đây là thời điểm phù hợp để khép lại hành trình của mình với \hll\ ạ.''

Cả bố tôi và bác Hijaki đều im lặng, chờ đợi em ấy nói tiếp. Em ấy hít một hơi, rồi mỉm cười nhẹ. ``Cháu muốn thử thách bản thân ở những lĩnh vực mới, không chỉ bó buộc trong khuôn khổ cũ. Và trong bối cảnh này\dots\ cháu vẫn tiếp tục làm VTuber, nhưng với một con đường riêng của mình.''

Tôi nhìn Coco một lúc, rồi hỏi: ``Vậy nên em mới chọn nhà Hikawa làm trại bản doanh của mình à?''

``Vâng!'' cô nàng gật đầu chắc nịch. ``Cháu muốn có một nơi vừa đủ điên rồ nhưng cũng đủ tin cậy để tiếp tục hành trình mới. Với lại, nếu đã rời khỏi \hll, thì cháu vẫn cần một chỗ ở, đúng chứ?''

Bố tôi và bác tôi trao đổi ánh mắt với nhau, rõ ràng là vẫn chưa tiêu hóa hết câu chuyện. Tôi cũng hiểu chứ - mọi chuyện xảy ra quá nhanh, quá đột ngột. Nhưng\dots

\dots vấn đề đằng sau đó lại là một câu chuyện khá khó nói.

Nói như thế khá là trừu tượng khó hiểu, nên để tôi bắt đầu giải thích một chút về mọi chuyện thực sự đã xảy ra với nàng Rồng này: ``Để cháu nói cho hai người ngắn gọn thế này ạ.''

Tôi khẽ hắng giọng rồi bắt đầu kể.

Kiryu Coco-chan là một trong những thành viên táo bạo nhất của công ty mà tôi từng biết. Nếu tôi phải mô tả em ấy bằng một câu, thì đó sẽ là: ``Một streamer không có giới hạn.''

Tôi không nói quá đâu, vì nếu bác Hijaki-san hay bố tôi thử mở kênh của em ấy ra mà xem, họ sẽ thấy vô số thứ mà một idol ``bình thường'', hay thậm chí là một người có đầu óc bình thường không đời nào dám động vào.

Chẳng hạn, Coco-chan có cả một chuyên mục buổi sáng tên là AsaCoco, nghe thì tưởng giống chương trình tin tức bình thường, nhưng nội dung của nó thì\dots\ không thể gọi là bình thường được. Ngoài việc đăng tin về các sự kiện trong công ty, Coco còn bày ra hàng loạt trò bá đạo như bán ``AsaCoco dạng thuốc'' - tôi biết, đây là một trò đùa - rồi tổ chức quảng cáo bậy bạ, bày trò thử thách không tưởng, hay thậm chí là hướng dẫn người ta\dots\ chửi thề bằng tiếng Anh. Đúng vậy, một thần tượng, một tấm gương mà mọi người muốn noi theo lại đi dạy khán giả của mình những câu chửi bằng giọng Mỹ chuẩn.

Bố tôi đến đây liền nhíu mày: ``K-Khoan, cái gì cơ? Dạy chửi thề á?''

Nàng Rồng tự hào chống nạnh, gật đầu đầy kiêu hãnh: ``Chuẩn đét luôn ạ! Cháu còn có cả một giáo trình hoàn chỉnh cho việc đó nữa đấy!''

Cả ba chúng tôi chỉ có thể toát mồ hôi lạnh với ``bảng thành tích tội lỗi'' đồ sộ của em ấy.

Bác Hijaki-san đưa tay lên bóp trán, thở hắt ra: ``\dots Tao cứ tưởng chỉ có cái công ty mày làm mới có vấn đề, hóa ra người trong đó còn lắm trò hơn tao nghĩ đấy, Kuon ạ.''

Tôi cười khổ, rồi tiếp tục kể lể: ``Đó là chưa kể đến những buổi stream mà nội dung của nó có thể khiến nhân viên kiểm duyệt phải khóc thét đấy ạ. Một số thì hơi nhạy cảm một chút, một số thì có vẻ như đang thách thức ranh giới của YouTube, còn một số thì mang tính ``thử thách đạo đức'' khi động vào những chủ đề mà chả ai dám chạm tới đó hai người. Và thế là, một trong những điều mà Coco-chan nổi tiếng nhất chính là\dots\ bị YouTube đánh sập kênh liên tục đấy ạ\dots''

Bố tôi nghe xong lời giải thích của tôi liền khoanh tay, lắc đầu: ``Chú chưa bao giờ nghe thấy một người nào lại tự hào vì bị sập kênh hết.''

Coco giơ tay lên làm biểu tượng chiến thắng: ``Đó là phong cách của cháu mà, chú ơi! Dám làm thì dám chịu! Nhưng mà khổ nỗi là bị sập nhiều quá thì cũng khổ lắm ạ\dots''

``Thế thì ngay từ đầu đừng có làm mấy nội dung này chứ\dots'' cả ba chúng tôi đồng thanh, mặt tỏ ra ngao ngán.

Tôi thở dài một hơi rồi tiếp tục: ``Vậy nên, hai người cũng có thể đoán ra là sớm muộn gì công ty cũng sẽ có vấn đề với phong cách của em ấy ạ.''

Hijaki-san nhướng mày: ``Để tao đoán\dots\ họ bắt đầu muốn siết chặt nó lại đúng không?''

Nàng Rồng bĩu môi: ``Vâng, họ bắt đầu đặt ra thêm đủ loại quy định mới, bắt cháu phải tự kiềm chế hơn. Nhưng mà cháu đâu phải kiểu người có thể kiềm chế được lâu đâu ạ.''

Tôi gật đầu, tiếp tục giải thích: ``Và đó chính là lý do Coco-chan bắt đầu cảm thấy không thoải mái với công ty nữa. Từ một người có thể tự do làm mọi thứ mình thích, em ấy dần dần bị bó buộc bởi những quy định ngày càng chặt chẽ hơn ạ. Đặc biệt là sau một số sự cố, công ty bắt đầu kiểm duyệt nội dung stream gắt gao hơn, khiến cho nhiều ý tưởng của em ấy bị cắt ngay từ trong trứng nước đấy ạ.''

Bố tôi nghe vậy, quay sang vị khách rồi chống cằm suy nghĩ: ``Ừm\dots\ Chú cũng không rành lắm về mấy chuyện của công ty bọn cháu, nhưng mà nếu họ đã có quy định thì mình phải tuân theo chứ? Chẳng phải khi ký hợp đồng, cháu cũng đã chấp nhận điều đó sao?''

Coco-chan xua tay, vội bao biện: ``Vâng, cháu biết mà! Nhưng vấn đề là cháu không thích bị bó buộc như thế, chú hiểu không? Cháu gia nhập \hll\ vì cháu nghĩ đây là một nơi cho phép mình làm mọi thứ mà mình muốn. Nhưng đến khi nhận ra là có giới hạn, thì cháu lại cảm thấy không còn thoải mái như trước nữa\dots''

``Thì tự do cũng phải trong khuôn khổ thôi chứ\dots'' tôi ngán ngẩm đáp lại. ``Làm quái gì có cái gọi là tự do hoàn toàn đâu\dots''

Bác Hijaki-san hừ nhẹ, rõ ràng cảm thấy không ổn với suy nghĩ có phần ngây thơ của nàng Rồng  này: ``Chà, đúng là có những người không hợp với công ty lớn. Kiểu như mày, thích tự do hơn là làm theo khuôn khổ, đúng không?''

Coco-chan gật đầu chắc nịch: ``Chính xác! Cháu không ghét công ty đâu, nhưng cháu cũng không muốn phải kiềm chế bản thân quá mức. Cháu thích làm những thứ mà mình cảm thấy vui, nhưng mà\dots\ nếu tiếp tục ở lại, cháu sẽ không thể làm được điều đó nữa.''

Tôi tiếp lời, kết thúc bài giải thíc: ``Vậy nên em ấy đã quyết định tốt nghiệp. Một phần vì không muốn bị giới hạn, một phần vì muốn thử thách bản thân ở một môi trường mới. Và bây giờ thì\dots'' tôi thở dài, nhìn sang nàng Rồng, ``\dots em ấy lại chọn nhà mình làm đại bản doanh mới ạ.''

Hai người lớn đồng loạt quay sang nhìn Coco đang phấn khởi với những ánh mắt đầy nghi hoặc. Cô nàng nở một nụ cười đầy tự hào, giơ tay lên chào kiểu quân đội: ``Vâng! Cháu quyết định sẽ tiếp tục làm VTuber, nhưng theo cách của riêng mình! Và nơi khởi đầu mới của cháu chính là ở đây, ký túc xá nhà Hikawa! Mong được hai người giúp đỡ ạ!''

Ryujin lắc đầu, thở dài: ``Chú không biết đây là vinh hạnh hay là điềm xấu nữa\dots''

Hijaki khoanh tay, cười khẩy: ``Tao chỉ biết một điều là từ nay, cái nhà này sẽ không bao giờ có một ngày yên bình.''

Tôi nhún vai cam chịu: ``Nó đã không yên bình từ khi Aloe-chan chuyển vào nhà mình rồi ạ\dots''

\begin{center}
    \colorbox{red}{\color{white} Phân đoạn dưới đây chứa yếu tố nhạy cảm. Các bạn nên cân nhắc thật kỹ trước khi đọc.}\\
    \colorbox{red}{\color{white} Các bạn có thể bỏ qua phần này nếu không rành về chính trị.}
\end{center}

Bất chợt, bố tôi lên tiếng: ``Mà bố thấy con có nói về `một số sự cố' khiến công ty phải siết chặt nội dung stream hơn á. Cụ thể là mấy vụ nào thế?''

Tôi nhìn sang Coco rồi hỏi: ``Em còn nhớ hai cái stream mà em với Haato-chan làm từ hồi tháng Chín năm ngoái không?''

Coco-chan nhăn trán nghĩ ngợi một lúc, cố nhớ lại một hồi rồi crồi mắt sáng lên như vừa nhớ ra: ``À, có phải là cái buổi mà hai chị em phân tích kênh của nhau ạ?''

Tất nhiên, tôi biết chắc rằng bố và bác Hijaki chưa hiểu ý nếu chỉ nói như vậy, nên tiếp tục giải thích: ``Em còn nhớ đoạn mà hai người phân tích lượng người xem dựa trên quốc gia không?''

Coco gật đầu cái rụp: ``À, đúng rồi ha! Hai chị em có nhắc tới Đài Loan mà. Nhưng sau đó lại thấy có một đống bình luận tiếng Trung\dots\ chửi tụi em! Kỳ lạ ghê luôn á.''

Tôi cười nhạt: ``\dots Nó chẳng kỳ lạ đâu. Con nghĩ chắc bố với bác cũng mường tượng được vì sao mọi chuyện lại xảy ra như thế đúng không ạ?''

Hai người lớn khẽ gật đầu, gương mặt trầm xuống đầy suy tư. Bác tôi thậm chí còn lắc đầu, giọng pha chút thất vọng: ``Mày nói vậy khác nào chọc giận người ta rồi còn gì.''

Bố tôi cũng gật gù: ``Chẳng trách mà bị dính chưởng.''

Coco-chan nhìn hai người lớn, mặt nghệch ra chưa hiểu chuyện, nhưng vẫn tiếp tục kể: ``Rồi sau đó, mấy hôm sau á, cháu với Haato-paisen bị đình chỉ ba tuần lận! Mà cháu còn không hiểu vì sao luôn á!''

Tôi chống tay lên trán, cố gắng giải thích đơn giản nhất có thể: ``Pha lỡ miệng đó chính là nguyên nhân đấy. Đang yên đang lành, tự dưng lại vô tình đụng chạm đến một vấn đề nhạy cảm. Người ta tức giận là phải.''

Cô nàng Rồng trố mắt nhìn tôi: ``Ể? Anh nói gì, em vẫn chưa hiểu? Em mới chỉ nói đến quốc gia Đài Loan thôi mà?''

Tôi lắc đầu: ``Bản chất vấn đề đó phức tạp hơn em nghĩ nhiều. Một số người, nhất là người Trung Quốc, rất cực đoan trong chuyện này. Chỉ cần ai đó lỡ miệng nói khác quan điểm của họ thôi là họ sẽ xem như xúc phạm đến lòng tự tôn dân tộc.''

Lúc này, cả nhà Hikawa đều gật gù đồng ý. Bố tôi khoanh tay, giọng nặng nề: ``Người Etsunan bọn chú cũng chả ưa nổi cái kiểu `ái quốc mù quáng' như vậy.''

Bác Hijaki-san bồi thêm: ``Bọn nó không cần biết mày cố tình hay vô ý, chỉ cần dính vào là coi như có tội. Chẳng lạ gì chuyện sau đó tụi nó làm loạn lên.''

Coco-chan im lặng, gương mặt lộ vẻ bất ngờ. Có vẻ đến giờ em ấy mới thật sự hiểu chuyện gì đã xảy ra. ``Vậy\dots\ đó là lý do cháu bị đình chỉ?''

Tôi gật đầu, rồi quyết định nói luôn phần còn lại, mà có thể nghe vô hại: ``Không chỉ vậy đâu. Chính cái chuyện tưởng vô hại đó đã khiến công ty phải giải tán nhánh CN - tức là nhánh Trung Quốc đấy.''

Nàng Rồng mở to mắt, sốc đến mức thân người hơi hướng về phía tôi: ``Ể!? Thật á? Em không biết luôn á!''

Tôi chỉ thở dài, khoát tay như muốn cho qua: ``Mà thôi, chuyện cũng đã qua rồi. Giờ em hiểu thì tốt. Quan trọng là sau này rút kinh nghiệm.''

Bố tôi khoanh tay lại, giọng có phần nghiêm túc hơn: ``Tóm lại thế này, cháu cứ tập trung làm tốt thứ gì mà cháu đang làm đi. Đừng có dây dưa vào mấy chuyện không liên quan.''

Tôi gật đầu, tóm gọn lại ý chính: ``Nói theo bố anh tức là `biết nhưng không nói'. Thậm chí không cần biết thì càng tốt. Vẹn cả đôi đường, không gây phiền hà cho ai.''

Đến đây, cô nàng Rồng đanh đá liền cúi đầu, thở ra một hơi dài, rồi chậm rãi nói: ``\dots Cháu cũng sẽ chú ý hơn đến phát ngôn của cháu sau này ạ.''

Nhà Hikawa chúng tôi cười mỉm khi Coco-chan đã bắt đầu trưởng thành hơn đôi chút: ``Vậy là ổn rồi đấy.''

\begin{center}
    \colorbox{red}{\color{white} Hết phân đoạn nhạy cảm.}
\end{center}

\ct

Chúng tôi tiếp tục nói chuyện về thủ tục nhận phòng sắp tới của Coco-chan. Đang nói chuyện bình thường, bố tôi bất chợt hỏi: ``Thế, điểm mạnh của cháu là gì?''

Nàng Rồng lập tức đáp ngay mà không cần suy nghĩ: ``So với mọi người ở \hll\ thì cháu có thể nói được tiếng Anh trôi chảy ạ!''

``Ơ, thế tao tưởng người Nhật kém tiếng Anh lắm mà?'' bác Hijaki-san khẽ trố mắt ngạc nhiên.

Tôi bật cười và nhún vai: ``Thực sự mà nói thì em ấy không hẳn là người Nhật đâu ạ. Ở nơi em ấy sống thì tiếng Anh mới là ngôn ngữ chính của em ấy.''

``Thế á?'' cả hai người nghiêng đầu hỏi với sự ngạc nhiên.

``Con có điêu đâu,'' tôi nhún vai. ``Em ấy còn làm hẳn một series về meme tiếng Anh, ngoài việc còn bày cho người khác chửi thề tiếng Anh nữa đấy.''

``Anh ấy nói đúng đó. Ở trong sơ yếu lý lịch của cháu cũng ghi như thế,'' Coco-chan cũng bổ sung thêm.

Bác Hijaki-san tỏ vẻ khá ngưỡng mộ với khả năng ngôn ngữ của em ấy, rồi bác hỏi tiếp: ``Vậy chắc mày học tiếng Nhật lâu lắm rồi nhỉ?''

Cô nàng chỉ khẽ gãi đầu, cười hơi ngượng ngùng: ``Cháu cũng học từ mấy năm trước rồi ạ\dots''

Tôi liền suồng sã dự đoán: ``Chắc em mê mẩn văn hóa Nhật, rồi sau đó quyết tâm học ngôn ngữ nước bạn, rồi sử dụng kỹ năng đó để xin vào làm việc ở \hll\ chứ gì?''

Coco-chan lập tức bật cười và trả lời một cách  ``Chuẩn quá! Sao anh biết?''

Tôi chỉ nhún vai: ``Thì anh là Tổng quản lý của mấy đứa mà, sao không biết. Tiếng Nhật cũng có phải là ngôn ngữ mẹ đẻ của anh đâu, nên anh hiểu chứ.''

Bố tôi tiếp tục hỏi một câu mà có lẽ ai cũng đang tò mò: ``Dù gì cháu cũng sắp nghỉ làm ở \hll\ rồi, vậy cháu định làm gì tiếp theo?''

Em ấy trả lời ngay lập tức, không một chút do dự: ``Như cháu nói rồi đó, cháu vẫn sẽ làm VTuber thôi ạ! Nhưng lần này cháu đang suy nghĩ với một vẻ ngoài khác\dots''

Tôi vô thức lẩm bẩm một câu theo bản năng: ``Với tướng của em và cái cách mà em `đột kích' cái nhà ban nãy, anh nghĩ em làm xã hội đen là hợp đấy\dots\ Thiếu mỗi cái chùy sắt thôi là đủ bộ\dots''

Không ngờ em ấy lại nghe thấy!

``Đó! Chính nó! Cảm ơn anh nha, Kuon-paisen!''

Tôi chết lặng một giây. Hình như tôi vừa vô tình hiến kế cho hình tượng ảo tiếp theo của em ấy à?

\dots Mà thôi kệ. Ít nhất em ấy cũng có việc làm là tốt rồi.

\ct

Sau một hồi trò chuyện và hoàn tất thủ tục bàn giao phòng ốc cho vị khách mới của ký túc xá, tôi dẫn Coco-chan lên căn phòng mà em ấy sẽ ở trong vài ngày tới.

Nó nằm ở tầng bảy, với thiết kế tương tự các tầng dưới - mỗi bên hành lang có một cầu thang. Căn phòng này có tông màu cam, đúng với màu chủ đề của nàng Rồng tóc cam này, và được trang bị những tiện ích cơ bản giống như phòng của tôi và bác Hijaki-san.

Dù vậy, căn phòng này vẫn còn trống vắng vài chỗ, vì em ấy vẫn chưa thực sự tốt nghiệp, nên đồ đạc em ấy chưa chuyển hết về đây.

Tôi chỉ về khoảng trống gần đệm ở góc phòng và nói: ``Đây sẽ là chỗ để PC của em. Phòng khá rộng nên em có thể để thêm vài ba thứ nữa.''

Coco-chan ngó nghiêng một vòng rồi lẩm bẩm: ``Em thấy phòng này cũng hơi hẹp nhỉ\dots''

Tôi nhướng mày nhìn em ấy, giọng có chút bất mãn: ``Anh thấy thế là rộng đấy. Phòng này còn rộng gấp đôi so với của anh mà, còn kêu ca gì nữa?''

``\dots Cũng phải.'' Cô nàng bật cười gãi đầu, có vẻ như cũng tự thấy mình hơi đòi hỏi quá. Nhưng rồi em ấy chợt nhớ ra một chuyện quan trọng hơn: ``Mà anh nói rằng em có thể gặp lại mọi người ở \hll\ đúng chứ?''

Tôi hơi ngập ngừng một chút, rồi trả lời một cách chắc chắn: ``Hiện tại thì chưa, nhưng chắc chắn là có thể. Một ngày nào đó. Theo lời YAGOO-san là vậy.''

Dừng một chút, tôi nói tiếp: ``Trước mắt thì em sẽ gặp một người bạn cùng ký túc với mình.''

Coco-chan nhướng mày tò mò: ``Ai vậy?''

Tôi nheo mắt cười nhẹ: ``Một người mà có thể em đã từng nghe qua hoặc từng thấy trong các meme Reddit. Thuộc thế hệ thứ Năm.''

Em ấy chớp mắt, rồi lập tức đoán ra ngay mà chẳng cần suy nghĩ: ``Mano Aloe-chan đúng không?''

Tôi hơi giật mình: ``\dots Nhanh vãi.''

Cô nàng bật cười thích thú: ``Em có làm một vài tập meme Reddit có mặt em ấy nên biết. Với cả Polka-chan cũng từng kể về em ấy suốt.''

``Vậy thì chắc anh nghĩ hai người có thể làm bạn với nhau nhanh đấy.''

Tôi bắt đầu rời khỏi phòng để hoàn thiện nốt hồ sơ nhận phòng, nhưng trước khi đi vẫn không quên dặn: ``Em cứ nhìn quanh phòng đi, thiếu hay hỏng cái gì thì cứ báo. Đến hôm em chuyển về chính thức thì gọi thợ sửa một thể.''

Coco-chan giơ tay chào, vui vẻ reo lên: ``Vâng\~{}''

\il

Vậy là ký túc xá \hll\ của gia đình Hikawa chúng tôi và Tanigo-san đã chính thức kết nạp thêm thành viên thứ hai.

Đó là một cô Rồng có tính cách bạo dạn, nhưng khá trẻ con Kiryu Coco.

Cơ mà\dots\ tôi cũng hy vọng sẽ không có thêm bất cứ ai phải chuyển về đây nữa.

Bác Hijaki-san nói rằng đây là ``đày'' thì cũng không sai, nhưng\dots\ chính sách công ty là thế mà, đành phải chịu vậy.

Sau Mano Aloe-chan, giờ đến lượt Coco-chan cũng trở thành một phần của gia đình Hikawa. Dĩ nhiên, về mặt pháp lý, em ấy chỉ là khách trọ tạm thời, nhưng tôi hiểu rõ một khi ai đã bước vào đây, thì khó mà thoát khỏi cái vòng dây nhân duyên này.

Nhớ lại ngày hồi nàng Succubus chuyển vào, tôi đã nghĩ đó chỉ là một trường hợp cá biệt, một người lạc lõng cần chỗ nương tựa sau những sóng gió. Nhưng bây giờ nhìn lại, tôi mới nhận ra mình đã sai. Cái ký túc xá này không đơn thuần là một chỗ ở tạm thời, mà đang dần trở thành một mái nhà thực sự cho những kẻ không còn nơi nào để về.

Và Coc-chano cũng vậy.

Có lẽ em ấy vẫn chưa nhận ra, nhưng ngay khoảnh khắc em ấy bước chân vào nơi này, em ấy đã không còn là một người bị đẩy ra ngoài nữa. Mà là một người được chào đón.

Tôi bật cười khẽ.

``Gia đình Hikawa à\dots''

Không ngờ mình lại có ngày mở rộng cái ``gia đình'' này theo cách chẳng giống ai.

%==========================================TRUYỆN PHỤ=========================================
\omk

Cuối cùng, ngày đó cũng đã tới.

Buổi lễ tốt nghiệp của Kiryu Coco-chan. Đây là một buổi lễ tốt nghiệp với đủ đầy đủ mọi cung bậc cảm xúc, từ hài hước đến xúc động.

Như bất kỳ buổi lễ tốt nghiệp nào, nó có những khoảnh khắc ôn lại kỷ niệm cũ, những câu chuyện vui và talk-show rôm rả. Dù Coco-chan chỉ gắn bó với mọi người hơn một năm rưỡi, nhưng quãng thời gian đó đủ dài để tạo ra vô số kỷ niệm khó quên. Họ cùng nhau xem lại những đoạn clip cũ, từ những khoảnh khắc ngớ ngẩn trên stream đến những lần nàng Rồng này gây bão với các meme \textit{Reddit Shitpost Review}, rồi cùng bật cười khi nhớ lại những pha trò lố của em ấy.

Tất nhiên, cũng không thiếu những ``tai nạn'' hài hước mang thương hiệu \hll. Như việc một thành viên bên nhánh tiếng Anh - gợi ý nhé, Cá mập-chan - đến muộn do\dots\ nhìn nhầm múi giờ. Hay có người đang phát biểu cảm nghĩ mà lại lỡ đà biến thành một pha tấu hài bất đắc dĩ.

Nhưng rồi, giữa những tiếng cười, không ai có thể ngăn nổi những giọt nước mắt khi từng người lần lượt nói lời chia tay với nàng Rồng. Đối với nhiều thành viên, Coco-chan không chỉ là một đồng nghiệp, mà còn là một người chị đáng tin cậy, là chỗ dựa tinh thần giữa vô vàn sóng gió. Kanata, người bạn cùng phòng đã gắn bó với em ấy từ ngày đầu tiên, khóc đến mức không thể nói nên lời. Towa cố tỏ ra mạnh mẽ, nhưng giọng nói vẫn nghẹn lại. Thậm chí cả những thành viên ít khi thể hiện cảm xúc nhất cũng không giấu nổi nỗi buồn.

Và điều khiến tôi bất ngờ nhất - chính là YAGOO.

Vị CEO lúc nào cũng tỏ ra điềm tĩnh ấy lại không thể kìm được nước mắt khi chứng kiến nàng Rồng rời đi. Tôi chưa từng thấy ông ấy xúc động đến vậy. Nhưng nghĩ lại thì cũng phải thôi. Em ấy không chỉ là một tài năng của công ty, mà còn là một chiến binh mạnh mẽ, là người đã mở ra những con đường mới mà chưa ai dám đi.

Buổi lễ khép lại với gần năm trăm nghìn người theo dõi trực tiếp - một kỷ lục trong thế giới VTuber vào thời điểm đó.

\ct

Đó là chuyện của ngày hôm qua. Còn bây giờ, tôi và nàng Rồng đang trên đường về Kawachi, mang theo hành lý và cả những cảm xúc chưa kịp nguôi ngoai.

Tôi đã để em ấy có thời gian gặp lại các thành viên Thế hệ thứ Tư lần cuối trước khi rời đi. Sau khi chia tay mọi người, tôi bảo Coco-chan hóa rồng và cùng tôi bay về trung tâm thành phố Kawachi, nơi em ấy sẽ bắt đầu một cuộc sống mới ở quê nhà.

Tôi liếc nhìn nàng Rồng bên cạnh, vẫn còn đôi chút trầm tư.

``Kana-chan chắc khóc nhiều lắm nhỉ?'' tôi mở lời.

Nàng Rồng khẽ gật đầu, giọng vang: ``Vâng\dots\ Cả Towa-chan nữa\dots\ Em sẽ nhớ họ lắm.''

``Thôi nào, không phải buồn. Em biết đây không phải là cuộc chia ly vĩnh cửu, đúng chứ? Rồi sẽ có ngày mọi người gặp lại nhau thôi.''

Nghe vậy, em ấy hít một hơi thật sâu rồi nở nụ cười kiên định. ``Anh nói cũng phải\dots\ Em phải làm gì đó để cho họ không cảm thấy thất vọng!''

``Tốt. Nhưng làm ơn đừng chơi kiểu sát nhân hay côn đồ nữa, nguy hiểm lắm.''

Coco-chan bĩu môi: ``Ể\~{}\dots Thế còn-''

Tôi liền ngắt giọng: ``Buôn ma túy cũng không được. Làm người tốt đi cho anh nhờ cái.''

``Đúng là kể cả rời công ty rồi vẫn bị anh kiểm soát kìa\dots\ Hông vui\dots\~{}''

Tôi bật cười. Ít nhất, em ấy vẫn là Coco-chan - cô nàng Rồng cam đầu đất mà tôi biết.

\ct

Chẳng mấy chốc mà chúng tôi đã bay về đến quận Kyuukami - phố mà tôi đang ở. Nhưng hình như Coco-chan vẫn chưa quen với việc hạ cánh khi bay với tốc độ cao, thế nên\dots

*RẦM!*

Em ấy đâm sầm vào tầng ba của ký túc xá nhà tôi, tạo nên một tiếng động lớn đến mức người đi đường bên dưới và các nhân viên cửa hàng gần đó đều giật mình hốt hoảng.

``\textit{Mình đáng ra phải tính đến chuyện này rồi mới phải\dots}'' tôi lác đầu, cố trấn tĩnh lại sau cú va chạm bất ngờ. May mà tầng ba đang không có ai ở (vì mẹ tôi có lẽ đang ở tầng năm, còn em tôi đang đi học), nên không có thương vong nào xảy ra.

Và quan trọng hơn, tôi không phải xây lại cả cái ký túc xá, hay toàn bộ cả căn nhà này.

Ngay khi Coco-chan biến trở lại thành người, một cơ chế đặc biệt mà tôi đã lắp đặt từ trước liền kích hoạt. Những bức tường bị sụp dần dần tái tạo lại theo cấu trúc ban đầu, từng viên gạch, từng mảnh kính vỡ đều tự động khôi phục. Đây chính là máy tái tạo công trình mà tôi đã nghiên cứu suốt mấy tuần qua, kể từ lúc Coco-chan xông vào nhà tôi ngày hôm đó. Lần đầu tiên thử nghiệm trong thực tế\dots\ và nó hoạt động hoàn hảo!

``\textit{Hệ thống này có vẻ ổn đấy. Có khi mình nên áp dụng vào văn phòng \hll\ nữa\dots}'' tôi lẩm bẩm. Chắc chắn tôi sẽ phải áp dụng nó vào công ty của chúng tôi, chứ không một tháng mà phải xây lại cả công ty tới bốn, năm lần thì có lấy hết lương của mọi người cũng chẳng xây kịp nữa!

Nhưng, chưa kịp tận hưởng thành quả thì\dots

``\textbf{\vie{Có chuyện gì mà ồn ào ngoài kia thế!?}}'' cả mẹ tôi, bác Hijaki-san và Aloe-chan cùng lao ra ngoài, mặt ai nấy đều hốt hoảng.

``Chuyện gì xảy ra vậy!?'' nàng Succubus hỏi, mắt đảo nhanh về phía bức tường vẫn còn đang dần khôi phục.

Từ trong đống gạch vụn đang được sửa chữa lại, Coco lảo đảo bước ra, vẫn còn hơi choáng váng: ``Chóng mặt quá\dots\ Mịa nó\dots''

Tôi khoanh tay nhìn em ấy: ``Đi đứng lần sau cẩn thận đi chứ.''

Bác tôi nhíu mày: ``Mày có sao không đấy!?''

Mẹ tôi lo lắng hỏi: ``Có ổn không cháu?''

``Dạ, cháu ổn ạ\dots\ Tường nhà ở đây cứng vãi đạn\dots'' Coco đáp, cố lấy lại cân bằng.

Trời đất, không ai hỏi han tình hình tôi sao? Nhưng mà thôi kệ, có lẽ họ cũng đã biết được cô nàng Rồng ấy sẽ chuyển về chung sống với chúng ta rồi.

Ngoại trừ\dots\ Aloe-chan. Vì con bé có mấy khi ở nhà đâu.

Hai người lớn nhanh chóng dìu em ấy vào nhà để nghỉ ngơi, còn tôi với Aloe-chan đứng lại lo chuyện hành lý. Khi thấy bóng dáng Coco-chan, nàng Succubus lập tức liếc mắt nhìn tôi đầy nghi hoặc: ``Kuon-kun, ai đây?''

Tôi nhấc một thùng đồ lên, bình thản đáp: ``Một thành viên đã tốt nghiệp nữa từ \hll, giờ chuyển đến ký túc xá nhà ta.''

Aloe nheo mắt quan sát kỹ hơn: ``Sao em có cảm giác quen quen nhỉ\dots\ Cặp sừng rồng này nhìn quen lắm\dots''

Tôi mỉm cười: ``Chắc em cũng từng nghe qua Kiryu Coco-chan rồi nhỉ?''

Đôi mắt nàng Succubus hơi mở to. ``Hình như mấy người bạn của ta từng kể về người đó rồi\dots\ Mà chị ấy trông như thế nào vậy, chủ nhà?''

Tôi nhún vai. ``Cứ hỏi em ấy khắc biết. Anh nghĩ hai người sẽ thân nhau đấy.''

Aloe-chan đeo tai nghe lên, ngẫm nghĩ một lát rồi định hỏi tiếp. Nhưng trước khi em ấy kịp mở lời, tôi đã nhanh chóng chặn lại: ``Mà nói đến khách thì\dots\ phụ anh xách đống đồ hộp này được không?''

``P-phải ha!'' cô nàng giật thót, vội vã đáp. ``Xếp đống đồ này lên thang máy rồi lên thôi!''

Đúng, chúng tôi còn lắp cả thang máy nữa. Chứ các bạn nghĩ chúng tôi sẽ leo thang bộ từ đây lên đến tận tầng bảy, và thậm chí là cao hơn nữa, nếu kịch bản xấu xảy đến ư?

Vậy là chúng tôi cùng nhau chuyển hết đồ của Coco lên tầng bảy, nơi em ấy sẽ ở từ nay về sau.

\ct

Suốt cả buổi tối hôm đó, Aloe-chan và Coco-chan trò chuyện khá rôm rả. Tôi ngồi một góc, vừa kiểm tra lại đống hành lý còn chưa sắp xếp xong, vừa lắng nghe hai người họ tám chuyện.

``Vậy là chị thật sự từng làm talk-show trên kênh chính của \hll\ luôn hả?" Aloe-chan hỏi, mắt sáng rỡ.

Nàng Rồng tiền bối tự hào đáp lại: ``Chuẩn luôn! \textit{Reddit Shitpost Review} là chương trình của chị đó! Toàn là meme xịn, cười đến muốn bể bụng luôn à!''

Nàng kusogaki nghiêng đầu: ``Cái đó là kiểu chương trình gì vậy? Kiểu kể chuyện cười à?''

``Không hẳn. Chủ yếu là chị đọc mấy cái meme mà fan làm về \hll, rồi phân tích chúng theo kiểu chuyên sâu-''

Tôi chen vào: ``Chuyên sâu? Theo nghĩa là\dots\ soi từng khung hình một à?''

``Hề hề\~{}'', cũng có thể lắm! Nhưng quan trọng là bọn chị tạo ra một không gian để cả công ty cùng cười với nhau. Nhờ thế mà mọi người đoàn kết hơn!'' em ấy vừa nói vừa vỗ ngực tự hào.

Thật sự thì chương trình đó làm cho tôi muốn đổ mồ hôi hột đấy\dots

Aloe-chan nhướng mày, nhìn em ấy với ánh mắt có chút nghi ngờ. ``Hmm\dots\ nghe giống chương trình tấu hài thì đúng hơn\dots''

Coco-chan bật cười. ``Chính xác! Mà nói mới nhớ, em có thấy cái meme nào của chính mình chưa?''

Nàng Succubus giật mình: ``Hả? Meme của ta á? Ta có nổi đến mức đó đâu\dots''

Nàng Rồng híp mắt: ``Em nghĩ thế thôi. Nhưng idol nào mà chả có meme. Để mai chị lục lại coi có cái nào của em không-''

``\textbf{Không cần đâu!!!}'' Aloe-chan hoảng hốt xua tay, mặt đỏ bừng vì ngượng ngùng, làm cho tiền bối phá lên cười.

Tôi chỉ ngồi đó, khoanh tay quan sát cách hai chị em trò chuyện với nhau. Xem ra hai người này hợp nhau nhanh hơn tôi tưởng.

Rồi, từ chủ đề meme, cuộc trò chuyện dần chuyển qua những kỷ niệm của Coco-chan khi còn làm ở \hll.

``Nhưng mà\dots'' Aloe-chan chống cằm, trầm ngâm: ``Rốt cuộc vì sao chị lại rời đi vậy?''

Bầu không khí bỗng chốc chùng xuống. Tiền bối của em ấy im lặng một lát, ánh mắt hơi dao động, nhưng rồi em ấy nở một nụ cười nhẹ: ``Cũng có nhiều lý do lắm\dots\ nhưng quan trọng nhất là chị muốn tạo ra một con đường mới cho chính mình.''

``Vậy à\dots''

Nhận thấy không khí có phần nặng nề, tôi đành lên tiếng phá tan bầu không khí trầm lắng: ``Thôi nào, đừng có mặt ủ mày chau thế. Giờ thì Coco đã là người của nhà Hikawa rồi, lo gì nữa?''

Coco bật cười, dường như hùa theo với lời của tôi: ``Ừ nhỉ! Giờ chị là thành viên của `ký túc xá bất đắc dĩ Hikawa' rồi mà!''

Em không cần phải nhắc cái chi tiết đó đâu\dots\ Đến anh cũng không biết tại sao nhà của anh lại biến thành cái nhà trọ nữa\dots

Aloe-chan nhếch mép, cười nửa miệng: ``Heh, thế thì ta sẽ theo dõi xem chị có sống sót qua được cái ký túc này không.''

Coco-chan khoanh tay, lên giọng như một nữ xã hội đen: ``Cô em định thách thức à? Được thôi, chị đợi đấy, nhóc Succubus!''

Tôi chống tay lên trán, thở dài ngán ngẩm: ``Rồi xong\dots\ lại thêm một cặp thích cà khịa nhau\dots''

Dù gì thì, chào mừng em đến với gia đình Hikawa, Kiryu Coco-chan!

Tôi thực sự không biết liệu gia đình tôi sẽ đón tiếp ai nữa hay không đây\dots

\end{document}